% part: intuitionistic-logic
% chapter: soundness-completeness
% section: truth-lemma

\documentclass[../../../include/open-logic-section]{subfiles}

\begin{document}

\olfileid[pt]{int}{sc}{tru}

\olsection{O lema da verdade}

O lema a seguir relaciona a satisfação no modelo canônico com quais
!!{formula}s são !!{element}s do conjunto primo $\Delta$.

\begin{lem}\ollabel{lem:truth}
  Se $\Delta$ é primo, então $\mSat{M(\Delta)}{!A}[\sigma]$ se, e somente se,
  $\Delta(\sigma) \Proves !A$.
\end{lem}

\begin{proof}
  Por indução sobre~$!A$.
  \begin{enumerate}
    \item \indcase{!A}{\lfalse}{Como $\Delta(\sigma)$ é primo, ele
      é consistente, logo $\Delta(\sigma) \Proves/ \indfrm$. Por
      definição, $\mSat/{M(\Delta)}{\indfrm}[\sigma]$.}
      \item \indcase{!A}{p}{Pela definição de $\mSat{{}}{}$,
        $\mSat{M(\Delta)}{\indfrm}[\sigma]$ se, e somente se, $\sigma \in V(p)$,
        O.s.\@, $\Delta(\sigma) \Proves \indfrm$.}
      \item \indcase!{!A}{\lnot !B}{}
      \item \indcase{!A}{!B \land
        !C}{$\mSat{M(\Delta)}{\indfrm}[\sigma]$ se, e somente se,
        $\mSat{M(\Delta)}{!B}[\sigma]$ e
        $\mSat{M(\Delta)}{!C}[\sigma]$. Pela hipótese de indução,
        $\mSat{M(\Delta)}{!B}[\sigma]$ se, e somente se, $\Delta(\sigma) \Proves
        !B$, e analogamente para~$!C$. Mas $\Delta(\sigma) \Proves !B$
        e $\Delta(\sigma) \Proves !C$ se, e somente se, $\Delta(\sigma)
        \Proves \indfrm$.}
      \item \indcase{!A}{!B \lor
        !C}{$\mSat{M(\Delta)}{\indfrm}[\sigma]$ se, e somente se,
        $\mSat{M(\Delta)}{!B}[\sigma]$ ou
        $\mSat{M(\Delta)}{!C}[\sigma]$. Pela hipótese de indução,
        isso vale se, e somente se, $\Delta(\sigma) \Proves !B$ ou
        $\Delta(\sigma) \Proves !C$. Devemos mostrar que isso, por
        sua vez, vale se, e somente se, $\Delta(\sigma) \Proves \indfrm$. A
        direção da esquerda para a direita é clara. A direção da direita para a esquerda
        decorre do fato de $\Delta(\sigma)$ ser primo.}
      \item \indcase{!A}{!B \lif !C}{Primeiro, a contrapositiva da
        direção da esquerda para a direita: suponha que $\Delta(\sigma) \Proves/ !B
        \lif !C$. Então também $\Delta(\sigma) \cup \{!B\} \Proves/
        !C$. Como $\tuple{!B, !C}$ é $\tuple{!B_n, !C_n}$ para
        algum~$n$, temos $\Delta(\sigma.n) = (\Delta(\sigma) \cup
        \{!B\})^*$, e $\Delta(\sigma.n) \Proves !B$, mas $\Delta(\sigma.n) \Proves/
        !C$. Pela hipótese de indução, $\mSat{M(\Delta)}{!B}[\sigma.n]$
        e $\mSat/{M(\Delta)}{!C}[\sigma.n]$. Como
        $R\sigma(\sigma.n)$, isso significa que
        $\mSat/{M(\Delta)}{\indfrm}[\sigma]$.

        Agora suponha que $\Delta(\sigma) \Proves !B \lif !C$ e seja
        $R\sigma\sigma'$. Como $\Delta(\sigma) \subseteq
        \Delta(\sigma')$, temos: se $\Delta(\sigma') \Proves
        !B$, então $\Delta(\sigma') \Proves !C$. Em outras palavras, para
        todo $\sigma'$ tal que $R\sigma\sigma'$, ou
        $\Delta(\sigma') \Proves/ !B$ ou $\Delta(\sigma') \Proves
        !C$. Pela hipótese de indução, isso significa que, sempre que
        $R\sigma\sigma'$, ou $\mSat/{M(\Delta)}{!B}[\sigma']$ ou
        $\mSat{M(\Delta)}{!C}[\sigma']$, O.s.\@,
        $\mSat{M(\Delta)}{\indfrm}[\sigma]$.}
  \end{enumerate}
\end{proof}

\end{document}
